% ==================
% Lernstrategien: Leitner-System & Spaced Repetition
% ==================

\setauthor{\secondauthor}

\section{Lernstrategien}
\label{chap:lernstrategien}
\setauthor{\secondauthor}

Im Rahmen der Entwicklung einer digitalen Lernplattform für HTL-Maturanten und HTL-Maturantinnen müssen Lernstrategien gewählt werden, die über reines Auswendiglernen hinausgehen. Die Matura zeichnet sich durch enormen Stoffumfang und hohe Komplexität aus. Gefragt ist nicht nur isoliertes Faktenwissen, sondern vor allem das Verständnis von Zusammenhängen und die Anwendung von Wissen auf neue, vernetzte Problemstellungen \cite{StudySmarterApp}.

\subsection{Spaced Repetition}

Spaced Repetition, auch verteiltes Lernen genannt, ist eine kognitionswissenschaftliche Methode, die Informationen durch zeitlich versetzte Wiederholungen im Langzeitgedächtnis verankert. Sie stellt die wissenschaftliche Antwort auf die Ebbinghaus’sche Vergessenskurve dar \cite{pfhSpacedRepetition}.

\textbf{Allgemeine Funktionsweise}

Die Ebbinghaus’sche Vergessenskurve (1885) beschreibt, dass neu gelernte Informationen exponentiell vergessen werden: Nach 20 Minuten bleiben nur 58\,\% erhalten, nach 1 Stunde 44\,\%, nach einem Tag ca. 33\,\%. \cite{StudySmarterSpaced}.

\begin{figure}[H]
    \centering
    \fbox{\includegraphics[width=0.7\textwidth]{SpacedRepetition.png}}
    \caption{Ebbinghaus’sche Vergessenskurve \cite{HeylamaSpacedRepetition}}
    \label{fig:SpacedRepetition}
\end{figure}

Diese Methode basiert auf dem Prinzip, dass Informationen genau dann wiederholt
werden, wenn das Gehirn kurz davor ist, sie zu vergessen. Ein praxisnahes Beispiel ist
das Lernen einer komplexen Formel aus der Nachrichtentechnik: Die erste Wiederholung
findet nach 24 Stunden statt, die zweite nach drei Tagen und die dritte nach einer
Woche. Mit jedem Erfolg vergrößert sich der zeitliche Abstand progressiv.\cite{pfhSpacedRepetition}

\textbf{Kontext der Matura}
Diese Methode erweist sich als extrem zeiteffizient, da sie das kurzfristige Auswendigler-
nen direkt vor der Prüfung durch eine nachhaltige Verankerung im Langzeitgedächtnis
ersetzt. Für die Maturavorbereitung ist dies aufgrund des langen Zeitraums essenziell,
um sicherzustellen, dass technische Grundlagen auch dann noch präsent sind, wenn am
Ende des Jahres hochkomplexe Anwendungsbeispiele gelöst werden müssen.\cite{StudySmarterSpaced}.

\subsection{Das Leitner-System}\cite{PocketPrepLeitner}

Das Leitner-System ist eine systematische Lernmethode, die die Organisation von Lern-
inhalten in verschiedenen Kategorien nutzt, um den Lernprozess zu strukturieren.

\textbf{Allgemeine Funktionsweise}

\begin{figure}[H]
    \centering
    \fbox{\includegraphics[width=0.7\textwidth]{LeitnerSystem.jpg}}
    \caption{Leitner-System mit 3 Boxen\cite{PocketPrepLeitner}}
    \label{fig:LeitnerSystem}
\end{figure}

Das System nutzt das Prinzip von physischen oder digitalen Lernkartei-Boxen, um den
Lernstoff nach dem individuellen Beherrschungsgrad zu sortieren. In einer digitalen
Anwendung wandern Karten bei einer richtigen Antwort in ein höheres Fach, das seltener
abgefragt wird, während sie bei Fehlern unmittelbar zurück in das erste Fach verschoben
werden\cite{hymerLeitner}.

\textbf{Kontext der Matura}
Obwohl das System eine gute Visualisierung des Fortschritts bietet, ist es in seiner klassischen Form oft zu starr. Für Maturanten und Maturantinnen es dennoch hilfreich, da die Einteilung
in Fächer den Fortschritt bei großen Themengebieten greifbar macht\cite{managementLeitner}.

\subsection{Microlearning}

Microlearning bezeichnet das Lernen in sehr kurzen, inhaltlich fokussierten Einheiten,
um umfangreiche Wissensbereiche handhabbar zu machen.

\textbf{Allgemeine Funktionsweise}
Microlearning zerlegt Stoff in kurze, fokussierte Einheiten (5–10 Minuten) für schnelles Feedback und niedrige Hemmschwellen \cite{SerloAbout}.

\textbf{Kontext der Matura}
Während dies den Einstieg erleichtert, ist die Methode für die Matura nur bedingt
geeignet. Die Gefahr besteht in einer Fragmentierung des Wissens: Die Matura verlangt
das Verständnis großer, vernetzter Systeme. Werden diese zu stark zerstückelt, verlieren
Schüler den Blick für die Kausalzusammenhänge. Microlearning versagt daher oft bei
der Vermittlung der für die Matura essenziellen Tiefe \cite{FeedbackTimeMultimediaLearning}.

\subsection{Kombination aus Leitner und Spaced Repetition}

Für den Maturatrainer wird eine Zusammensetzung aus dem Leitner-System und Spaced
Repetition implementiert. Diese Kombination löst die Schwächen der Einzelmethoden
auf:

\begin{itemize}
    \item \textbf{Struktur durch Leitner:} Die Einteilung in fünf digitale Fächer dient als intuitive
Benutzeroberfläche. Der/ die SchülerIn sieht jederzeit, wie viele Karten sich im Bereich
Beherrscht befinden, was die Selbstwirksamkeit stärkt\cite{LeitnerWikipedia}.
    \item \textbf{Effizienz durch Spaced Repetition:} Durch die Einteilung in die Leitner-
Fächer wird die Spaced Repetition automatisch gesteuert. Fragen, die falsch
beantwortet wurden, bleiben im ersten Fach, die häufig richtig beantworteten
Karten wandern in die höheren Fächer, wodurch die Intervalle für die Wiederholung
automatisch angepasst werden. Je höher das Fach, desto seltener die Wiederholung,
was die Effizienz maximiert\cite{hymerLeitner}.
\end{itemize}

Dieser Mix ist für die Matura ideal, er bietet die motivierende Struktur des Leitner-
Systems, kombiniert mit der Effizienz der Spaced Repetition, um sicherzustellen, dass
das Wissen nachhaltig verinnerlicht wird.
